\chapter{Proposta de desenvolvimento}

A proposta de desenvolvimento da solução busca conciliar modularidade, segurança, escalabilidade e capacidade de evolução. Para isso, será estruturada segundo uma arquitetura de microsserviços orientada a eventos, utilizando o \textit{Apache Kafka} como barramento de comunicação assíncrona entre os principais serviços. Essa abordagem permite reduzir o acoplamento entre os componentes, distribuir responsabilidades de domínio e processar tarefas de forma independente, favorecendo a resiliência a picos de carga e a escalabilidade horizontal.

Segundo Newman (2022), microsserviços constituem uma abordagem na qual uma aplicação é organizada como um conjunto de serviços pequenos e independentes, responsáveis por capacidades específicas e passíveis de implantação e evolução separadas. No EIGAM, essa organização é utilizada para separar responsabilidades relacionadas à interação conversacional, processamento financeiro, notificações, auditoria, integração com serviços externos e processamento assíncrono.

A arquitetura também será concebida de forma desacoplada das plataformas de mensageria. O \textit{WhatsApp} será utilizado como canal inicial por meio da \textit{Meta Cloud API}, porém as regras de negócio não dependerão diretamente dos formatos e recursos específicos desse provedor. Uma camada de adaptação será responsável por converter mensagens de diferentes canais para representações internas padronizadas e, no sentido inverso, adaptar as respostas produzidas pelo EIGAM às capacidades do canal utilizado. Dessa forma, novos serviços de mensageria poderão ser incorporados sem alterações significativas no núcleo funcional da aplicação.

A camada de inteligência artificial seguirá a mesma premissa de separação de responsabilidades. O modelo de linguagem não será tratado como fonte de verdade dos dados financeiros, mas como um recurso para interpretação de linguagem natural, classificação de intenção, recuperação contextual e geração textual. Cálculos financeiros, validações de autorização, mascaramento de valores e parte das análises preditivas permanecerão sob responsabilidade de mecanismos determinísticos ou modelos específicos. O \textit{Chat Service}, o \textit{Finance Service} e o \textit{Notification Service} utilizarão técnicas distintas conforme suas responsabilidades, mantendo a inteligência distribuída entre os serviços em vez de concentrá-la em um único componente.

O processo de desenvolvimento seguirá práticas ágeis e utilizará \textit{Test-Driven Development} (TDD) nas funcionalidades em que sua aplicação for adequada. Além dos testes tradicionais de software, serão previstos testes específicos para integrações entre microsserviços, adaptadores de mensageria, mecanismos de RAG, classificação de intenções, controles da camada conversacional e modelos preditivos.

\section{Tecnologias que serão utilizadas}

O Quadro~\ref{frame:tecnologias} apresenta as principais tecnologias previstas para o desenvolvimento do EIGAM. A seleção considera as necessidades funcionais e arquiteturais definidas nas etapas de modelagem, incluindo comunicação assíncrona, persistência heterogênea, integração multicanal, recuperação vetorial, análise preditiva, segurança e observabilidade.

Ressalta-se que serviços externos e modelos de inteligência artificial podem sofrer alterações de versão, disponibilidade ou condições comerciais ao longo do desenvolvimento. Dessa forma, sempre que possível, a arquitetura deverá utilizar interfaces internas que reduzam o acoplamento com fornecedores específicos e permitam substituições sem impacto significativo nas regras de negócio.

{\small
\begin{longtable}{|L{3.2cm}|L{2.8cm}|L{6.7cm}|L{1.5cm}|}
\caption{Tecnologias Utilizadas} \label{frame:tecnologias} \\
\continuesmark
\hline
\textbf{Nome} & \textbf{Fornecedor} & \textbf{Descrição} & \textbf{Versão} \\ \hline
\endfirsthead
\hline
\textbf{Nome} & \textbf{Fornecedor} & \textbf{Descrição} & \textbf{Versão} \\ \hline
\endhead
\endfoot
\sourcerow{4} \\
\endlastfoot

Git & Open Source & Sistema de controle de versão utilizado para rastrear alterações no código-fonte. & 2.x \\ \hline
GitHub & Microsoft & Plataforma de hospedagem de código-fonte, colaboração e automação de CI/CD por meio do GitHub Actions. & -- \\ \hline
Node.js & OpenJS Foundation & Ambiente de execução JavaScript utilizado nos serviços de integração, comunicação, orquestração e regras de negócio desenvolvidos em TypeScript. & 20 LTS \\ \hline
TypeScript & Microsoft & Superset de JavaScript com tipagem estática, utilizado para aumentar a segurança de tipos, a legibilidade e a manutenibilidade dos serviços. & 5.x \\ \hline
Fastify & Open Source & Framework HTTP para Node.js utilizado na implementação do API Gateway e das interfaces HTTP dos microsserviços. & 5.x \\ \hline
Python & Python Software Foundation & Linguagem utilizada nos componentes especializados de análise preditiva e séries temporais quando houver necessidade de bibliotecas específicas do ecossistema científico. & 3.x \\ \hline
Prophet & Open Source & Biblioteca utilizada na modelagem de séries temporais para usuários com histórico suficiente, permitindo considerar tendência, sazonalidade e intervalos de previsão. & -- \\ \hline
Ruptures & Open Source & Biblioteca destinada à detecção de mudanças estruturais em séries temporais, utilizada na implementação do algoritmo PELT. & -- \\ \hline
Apache Kafka & Apache / Upstash & Barramento de eventos distribuído responsável pela comunicação assíncrona entre os microsserviços. & -- \\ \hline
KafkaJS & Open Source & Cliente Kafka para Node.js utilizado pelos produtores e consumidores de eventos dos serviços implementados em TypeScript. & 2.x \\ \hline
Redis & Upstash & Armazenamento em memória utilizado para cache, sessões, tokens temporários e mecanismos de limitação de requisições. & 7.x \\ \hline
BullMQ & Open Source & Biblioteca de filas baseada em Redis utilizada no agendamento e processamento de sincronizações, notificações e tarefas periódicas. & 5.x \\ \hline
PostgreSQL & Open Source & Banco de dados relacional utilizado para informações estruturadas, usuários, consentimentos, auditoria e demais entidades relacionais. & 16 \\ \hline
pgcrypto & Open Source & Extensão do PostgreSQL utilizada como apoio à proteção de campos que necessitam de tratamento criptográfico no banco relacional. & built-in \\ \hline
pgvector & Open Source & Extensão do PostgreSQL utilizada no armazenamento, indexação e busca por similaridade de \textit{embeddings} empregados no pipeline de RAG. & 0.7+ \\ \hline
MongoDB & MongoDB, Inc. & Banco de dados documental utilizado para informações cujo formato se beneficia de maior flexibilidade estrutural, como transações e histórico conversacional. & 7.x \\ \hline
Drizzle ORM & Open Source & ORM TypeScript utilizado nas operações relacionais com PostgreSQL e no acesso tipado às entidades da aplicação. & -- \\ \hline
Mongoose & Open Source & ODM utilizado pelos serviços Node.js no acesso às coleções armazenadas no MongoDB. & 8.x \\ \hline
OpenAI API & OpenAI & API utilizada para acesso aos modelos de linguagem e geração de \textit{embeddings}. Inicialmente, prevê-se o uso de GPT-4o-mini no pipeline conversacional e text-embedding-3-small no mecanismo de RAG. & -- \\ \hline
Meta Cloud API & Meta & Integração inicial com a plataforma WhatsApp Business. Será encapsulada por um adaptador de canal para impedir dependência direta das regras de negócio em relação à API da Meta. & v21+ \\ \hline
Pluggy & Pluggy & API utilizada como intermediária para acesso aos dados compartilhados pelo usuário no ecossistema de \textit{Open Finance}. & v1 \\ \hline
Chart.js + @napi-rs/canvas & Open Source & Bibliotecas utilizadas para geração de gráficos em formato de imagem, possibilitando seu envio por canais de mensageria que suportem conteúdo multimídia. & 4.x \\ \hline
Docker + Docker Compose & Docker, Inc. & Plataforma de conteinerização utilizada para padronizar os ambientes de desenvolvimento, testes e implantação dos serviços. & 27+ \\ \hline
Nginx & Open Source & \textit{Reverse proxy} responsável pela terminação TLS e pelo encaminhamento de requisições externas ao API Gateway. & -- \\ \hline
Let's Encrypt & ISRG & Autoridade certificadora utilizada para emissão e renovação automatizada de certificados TLS. & -- \\ \hline
Helmet.js & Open Source & Conjunto de \textit{middlewares} utilizado para aplicação de cabeçalhos de segurança nas interfaces HTTP dos serviços Node.js. & 8.x \\ \hline
Zod & Open Source & Biblioteca de validação de esquemas utilizada para validar requisições, configurações e contratos de eventos. & 3.x \\ \hline
Vitest & Open Source & Framework de testes utilizado nos serviços TypeScript para testes unitários e parte dos testes de integração. & -- \\ \hline
Pytest & Open Source & Framework utilizado nos componentes Python para testes dos métodos estatísticos e modelos preditivos. & -- \\ \hline
Grafana Cloud & Grafana Labs & Plataforma utilizada para visualização de métricas, acompanhamento operacional e configuração de alertas. & -- \\ \hline

\end{longtable}
}

\subsection{Estratégia de Integração Multicanal}

A comunicação com plataformas de mensageria será implementada por meio de adaptadores responsáveis por isolar particularidades de cada provedor. O núcleo do EIGAM trabalhará com estruturas internas padronizadas para representar mensagens recebidas e mensagens a serem enviadas. Dessa forma, informações específicas do \textit{WhatsApp}, como identificadores, formatos de \textit{webhook}, recursos interativos e regras próprias da \textit{Meta Cloud API}, permanecerão concentradas no adaptador correspondente.

No fluxo de entrada, o adaptador converte o evento recebido do provedor para o formato interno utilizado pelo sistema e o encaminha ao processamento da aplicação. No fluxo de saída, a resposta produzida pelo núcleo é transformada pelo adaptador para o formato exigido pelo canal. A camada também deverá identificar as capacidades suportadas por cada plataforma, como texto, imagens, botões, formulários ou outros recursos interativos, permitindo aplicar alternativas quando determinada funcionalidade não estiver disponível.

A implementação inicial utilizará o \textit{WhatsApp}, porém a inclusão de um novo canal deverá exigir predominantemente a criação de um novo adaptador e sua configuração na camada de mensageria. Com isso, o \textit{Chat Service}, o \textit{Finance Service}, o \textit{Notification Service} e os demais componentes de domínio permanecerão independentes da plataforma utilizada pelo usuário.

\subsection{Estratégia de Implementação da Camada de Inteligência Artificial}

A implementação da camada de inteligência artificial seguirá a distribuição de responsabilidades definida na modelagem do sistema. O \textit{Chat Service} utilizará o modelo de linguagem para classificação semântica de intenções, orquestração do mecanismo de RAG e geração da resposta conversacional. O \textit{Finance Service} será responsável pela classificação e análise das transações, cálculo das projeções e detecção de comportamentos atípicos, enquanto o \textit{Notification Service} utilizará os resultados produzidos pelos demais componentes para definir e formatar comunicações proativas.

O pipeline de RAG utilizará \textit{embeddings} armazenados no PostgreSQL por meio do \textit{pgvector}. Os conteúdos utilizados para indexação serão produzidos a partir de dados normalizados e \textit{templates} internos da aplicação. Durante a consulta, os elementos semanticamente relevantes serão recuperados e poderão ser combinados com informações obtidas por consultas estruturadas antes da geração da resposta. Essa estratégia busca manter valores financeiros e resultados objetivos associados às fontes de dados da aplicação, utilizando o modelo de linguagem principalmente na interpretação e na apresentação das informações.

As análises preditivas não dependerão exclusivamente do LLM. Para usuários com histórico limitado, será utilizada uma estratégia de linha de base baseada em média móvel exponencialmente ponderada (EWMA). Quando houver histórico suficiente, poderão ser aplicados modelos de séries temporais, como o Prophet, para identificação de tendência e sazonalidade. A detecção de anomalias utilizará mecanismos estatísticos, como Z-score sazonal, e poderá utilizar o algoritmo PELT para identificar mudanças estruturais persistentes no comportamento financeiro.

Antes de uma resposta conversacional ser entregue ao usuário, serão executados controles determinísticos relacionados à autorização, exposição de valores monetários, consistência das informações e escopo da solicitação. O mascaramento de valores não autorizados não dependerá do modelo de linguagem. Respostas que contenham valores sem origem identificável nos dados estruturados ou no contexto fornecido ao modelo deverão ser rejeitadas ou regeneradas. Essa separação busca reduzir o risco de que uma resposta gerativa seja tratada como fonte autônoma de informação financeira.

\section{Metodologias e padrões de software}

Os padrões e práticas de engenharia de software são fundamentais para garantir a coerência, a manutenibilidade e a capacidade de evolução de sistemas distribuídos. Segundo Bass, Clements e Kazman (2022), decisões arquiteturais influenciam diretamente atributos como desempenho, segurança, disponibilidade e modificabilidade e, portanto, devem ser explicitamente definidas e documentadas. O Quadro~\ref{frame:padroes} apresenta as principais abordagens adotadas no EIGAM.

{\small
\begin{longtable}{|L{4.5cm}|L{10cm}|}
\caption{Metodologias e Padrões Adotados} \label{frame:padroes} \\
\continuesmark
\hline
\textbf{Nome} & \textbf{Descrição} \\ \hline
\endfirsthead
\hline
\textbf{Nome} & \textbf{Descrição} \\ \hline
\endhead
\endfoot
\sourcerow{2} \\
\endlastfoot

Metodologia Scrum & Framework ágil utilizado para organizar o desenvolvimento incremental em ciclos de trabalho, permitindo planejamento, acompanhamento, revisão e adaptação frequentes. \\ \hline
Arquitetura de Microsserviços Orientada a Eventos & Os principais serviços do EIGAM são separados por responsabilidades de domínio e utilizam eventos para os fluxos assíncronos. O Kafka atua como barramento central, enquanto comunicações síncronas pontuais podem utilizar APIs HTTP quando houver necessidade de resposta imediata. \\ \hline
Padrão MSC (\textit{Model-Service-Controller}) & Separação interna das responsabilidades em \textit{Controller}, responsável pela entrada e validação; \textit{Service}, responsável pela lógica de aplicação; e \textit{Model}, responsável pela representação e pelo acesso aos dados. \\ \hline
REST & Estilo arquitetural empregado em interfaces HTTP, integrações administrativas e comunicações síncronas pontuais entre componentes. \\ \hline
Adapter & Padrão utilizado na camada de mensageria para converter formatos específicos de cada provedor em contratos internos padronizados. Permite incorporar novos canais sem alterar as regras de negócio do EIGAM. \\ \hline
Strategy & Utilizado em pontos nos quais diferentes estratégias podem ser selecionadas conforme o contexto, como a escolha entre métodos de previsão para usuários com diferentes quantidades de histórico financeiro. \\ \hline
RAG (\textit{Retrieval-Augmented Generation}) & Abordagem que combina recuperação de informações relevantes em base vetorial e dados estruturados com a geração em linguagem natural. Seu objetivo é fornecer ao modelo contexto relacionado aos dados efetivamente disponíveis para o usuário. \\ \hline
Circuit Breaker e Retry & Mecanismos utilizados nas integrações externas para interromper temporariamente chamadas a serviços com falhas recorrentes e realizar novas tentativas controladas, reduzindo propagação de indisponibilidade. \\ \hline
TDD (\textit{Test-Driven Development}) & Prática na qual testes automatizados orientam a implementação de unidades de software. O ciclo consiste em criar um teste que inicialmente falha, implementar o comportamento necessário e posteriormente refatorar o código preservando os testes. \\ \hline

\end{longtable}
}

\section{Metodologia para testes de software}

A estratégia de testes do EIGAM será composta por diferentes níveis, adequados à natureza distribuída da aplicação. O TDD será utilizado principalmente nas regras de negócio e componentes cujo comportamento possa ser especificado de forma determinística. Segundo Richardson (2019), em arquiteturas de microsserviços, a validação isolada dos serviços deve ser complementada por testes das integrações, pois falhas podem ocorrer não apenas na lógica interna, mas também nos contratos estabelecidos entre componentes independentes.

Os testes unitários validarão regras de negócio, transformações de dados, autenticação, autorização, mascaramento de informações, adaptadores de mensageria e demais funções isoladas. Nos serviços TypeScript, esses testes serão implementados com ferramentas como o Vitest; componentes especializados implementados em Python utilizarão testes equivalentes com Pytest.

Também serão realizados testes de integração para verificar a produção e o consumo de eventos Kafka, persistência no PostgreSQL e no MongoDB, utilização do Redis, execução de tarefas agendadas pelo BullMQ e comunicação com serviços externos por meio de ambientes de teste ou mecanismos simulados. Conforme recomendado por Newman (2022), os contratos entre produtores e consumidores devem ser validados para reduzir o risco de que alterações em um serviço causem incompatibilidades silenciosas em outros componentes.

A camada multicanal deverá possuir testes de contrato para cada adaptador. Esses testes verificarão se mensagens recebidas de diferentes provedores são convertidas corretamente para o formato interno e se as respostas produzidas pelo EIGAM podem ser adaptadas aos recursos disponíveis em cada plataforma. O adaptador do \textit{WhatsApp}, utilizado no MVP, também deverá ser testado quanto ao processamento de \textit{webhooks}, envio de mensagens e tratamento de erros retornados pela \textit{Meta Cloud API}.

Os componentes de inteligência artificial exigem uma estratégia de validação adicional. A classificação de intenção será avaliada a partir de um conjunto de mensagens previamente rotuladas, permitindo acompanhar métricas como acurácia, precisão, revocação e F1 por intenção. Operações críticas relacionadas aos direitos do titular deverão manter identificação determinística e testes específicos que comprovem que tais comandos não dependem exclusivamente da classificação probabilística do LLM.

O mecanismo de RAG será avaliado separadamente da geração textual. Os testes deverão verificar se documentos semanticamente relevantes são recuperados para consultas conhecidas, se o limiar de similaridade é respeitado e se o comportamento de \textit{fallback} para dados estruturados ocorre quando não houver contexto vetorial suficiente. Também serão avaliados cenários em que o modelo tente gerar valores não presentes no contexto, garantindo que os mecanismos de validação impeçam a apresentação desses valores ao usuário.

Para os componentes preditivos, serão utilizados testes retrospectivos sobre séries históricas, comparando previsões com observações conhecidas. Métricas de erro adequadas ao comportamento dos dados, como MAE e RMSE, poderão ser acompanhadas para comparar as estratégias disponíveis e calibrar parâmetros. Os mecanismos de detecção de anomalias serão testados com séries contendo comportamentos conhecidos, permitindo avaliar a capacidade de identificar ocorrências atípicas sem produzir quantidade excessiva de falsos alertas.

Por fim, serão realizados testes direcionados aos mecanismos de segurança e controle da camada conversacional. Entre os cenários estão sessão expirada, tentativa de exibição de valores sem autorização, solicitação fora do escopo, tentativa de \textit{prompt injection}, resposta com informação monetária sem origem verificável e falhas de provedores externos. O objetivo é validar não apenas a resposta produzida em condições normais, mas também o comportamento seguro do sistema diante de situações adversas.

\section{Métricas de qualidade}

Segundo Bass, Clements e Kazman (2022), atributos de qualidade são propriedades mensuráveis de um sistema que indicam quão bem ele satisfaz as necessidades de seus \textit{stakeholders}. Em sistemas distribuídos orientados a eventos, características como disponibilidade, desempenho, segurança, modificabilidade e observabilidade devem ser consideradas desde a arquitetura. No EIGAM, também é necessário avaliar atributos específicos da camada de inteligência artificial e da integração multicanal. O Quadro~\ref{frame:metricas} apresenta os principais objetivos de qualidade adotados.

{\small
\begin{longtable}{|L{4cm}|L{10.5cm}|}
\caption{Métricas de Qualidade} \label{frame:metricas} \\
\continuesmark
\hline
\textbf{Fator de Qualidade} & \textbf{Objetivo} \\ \hline
\endfirsthead
\hline
\textbf{Fator de Qualidade} & \textbf{Objetivo} \\ \hline
\endhead
\endfoot
\sourcerow{2} \\
\endlastfoot

Manutenibilidade &
\begin{minipage}[t]{\linewidth}
\begin{itemize}[noitemsep,topsep=2pt,leftmargin=*]
    \item Organização em serviços com responsabilidades claramente delimitadas e separação interna por camadas.
    \item Utilização de TypeScript com tipagem estrita e validação de contratos por meio de schemas.
    \item Contratos de eventos versionados para reduzir incompatibilidades entre produtores e consumidores.
    \item Adaptadores de mensageria isolados das regras de negócio, permitindo inclusão ou substituição de canais com impacto reduzido no núcleo da aplicação.
\end{itemize}
\end{minipage} \\ \hline

Desempenho &
\begin{minipage}[t]{\linewidth}
\begin{itemize}[noitemsep,topsep=2pt,leftmargin=*]
    \item O recebimento de \textit{webhooks} dos canais deve ser confirmado rapidamente, encaminhando o processamento pesado para fluxos assíncronos sempre que possível.
    \item Como meta inicial, o pipeline conversacional completo, incluindo recuperação de contexto e chamada ao modelo de linguagem, deverá ser concluído em até 30 segundos em condições normais de operação.
    \item Processamentos em lote, como resumos periódicos, deverão ser distribuídos de forma a respeitar a janela de envio configurada sem comprometer consultas interativas.
\end{itemize}
\end{minipage} \\ \hline

Disponibilidade e Resiliência &
\begin{minipage}[t]{\linewidth}
\begin{itemize}[noitemsep,topsep=2pt,leftmargin=*]
    \item Utilização de \textit{circuit breakers}, políticas de repetição com \textit{backoff} e tratamento de indisponibilidade nas integrações externas.
    \item Filas de mensagens com falha para permitir análise e reprocessamento de eventos que não puderem ser concluídos automaticamente.
    \item \textit{Healthchecks} dos serviços conteinerizados e reinicialização controlada em caso de falha.
    \item Falhas em um provedor de mensageria não devem tornar indisponíveis os serviços financeiros internos nem impedir o funcionamento de outros canais futuramente integrados.
\end{itemize}
\end{minipage} \\ \hline

Segurança e Conformidade LGPD &
\begin{minipage}[t]{\linewidth}
\begin{itemize}[noitemsep,topsep=2pt,leftmargin=*]
    \item Proteção dos dados pessoais e financeiros em repouso e em trânsito por mecanismos criptográficos adequados.
    \item Token de privilégio com TTL de 1 hora para autorização temporária da visualização de valores monetários.
    \item Mascaramento determinístico de valores financeiros quando não houver autorização válida, independentemente da resposta produzida pelo LLM.
    \item Registro de operações relevantes em \textit{audit log} e suporte aos fluxos relacionados aos direitos do titular previstos na LGPD.
\end{itemize}
\end{minipage} \\ \hline

Portabilidade e Interoperabilidade &
\begin{minipage}[t]{\linewidth}
\begin{itemize}[noitemsep,topsep=2pt,leftmargin=*]
    \item Conteinerização dos serviços para reduzir diferenças entre os ambientes de desenvolvimento, testes e produção.
    \item Configuração externalizada por variáveis de ambiente e ausência de dependências de infraestrutura codificadas diretamente nas regras de negócio.
    \item Utilização de contratos internos padronizados para entrada e saída de mensagens, permitindo integração com diferentes plataformas de mensageria.
\end{itemize}
\end{minipage} \\ \hline

Usabilidade &
\begin{minipage}[t]{\linewidth}
\begin{itemize}[noitemsep,topsep=2pt,leftmargin=*]
    \item Interação por canais de mensageria já incorporados à rotina dos usuários, utilizando inicialmente o \textit{WhatsApp}.
    \item Respostas claras, objetivas e adaptadas às capacidades do canal utilizado, incluindo representações visuais quando suportadas e relevantes.
    \item Operações de consulta, gerenciamento de metas, preferências e direitos do titular acessíveis por interação conversacional ou recursos estruturados disponibilizados pelo canal.
\end{itemize}
\end{minipage} \\ \hline

Confiabilidade da Inteligência Artificial &
\begin{minipage}[t]{\linewidth}
\begin{itemize}[noitemsep,topsep=2pt,leftmargin=*]
    \item Avaliação periódica da classificação de intenção utilizando métricas como precisão, revocação e F1 por classe.
    \item Validação do mecanismo de RAG separadamente da geração de resposta, verificando a recuperação de contexto relevante e o comportamento quando não houver similaridade suficiente.
    \item Valores financeiros apresentados em respostas conversacionais devem possuir origem verificável nos dados estruturados ou no contexto disponibilizado ao modelo.
    \item Modelos preditivos devem ser avaliados por meio de testes retrospectivos e métricas de erro antes da utilização dos resultados em alertas ou recomendações.
    \item Controles de segurança da resposta devem funcionar independentemente do LLM, especialmente autenticação, autorização, mascaramento e comandos críticos relacionados aos dados do titular.
\end{itemize}
\end{minipage} \\ \hline

Observabilidade &
\begin{minipage}[t]{\linewidth}
\begin{itemize}[noitemsep,topsep=2pt,leftmargin=*]
    \item Coleta de métricas operacionais dos serviços e visualização centralizada por meio do Grafana Cloud.
    \item Logs estruturados com identificadores de correlação para permitir rastrear fluxos assíncronos entre os microsserviços.
    \item Alertas para crescimento da taxa de erro, aumento de latência, falhas em tarefas agendadas e indisponibilidade das integrações externas.
    \item Monitoramento de métricas próprias da camada de IA, como distribuição das intenções, solicitações de esclarecimento, falhas de recuperação de contexto e acionamentos de controles de segurança.
\end{itemize}
\end{minipage} \\ \hline

\end{longtable}
}
